\section*{Chapitre \ref{ch:b1:heat-budget} --- Le bilan thermique de la Terre}

\begin{solution}{exo:b1:heat-budget:1}
$q = k\,\mathrm dT/\mathrm dz = 2.5\times\qty{0.022}{K.m^{-1}} =
\qty{0.055}{W.m^{-2}} = \qty{55}{mW.m^{-2}}$. Through \qty{1}{km^2} $=
\qty{e6}{m^2}$~: \qty{5.5e4}{W}, soit \qty{55}{kW}.
\end{solution}

\begin{solution}{exo:b1:heat-budget:2}
Continents $\num{2.073e14}\times0.0709 = \qty{1.470e13}{W} = \qty{14.70}{TW}$~;
océans $\num{3.028e14}\times0.1054 = \qty{31.92}{TW}$~; total \qty{46.61}{TW}, et
le flux moyen $\num{4.661e13}/\num{5.101e14} = \qty{91.4}{mW.m^{-2}}$ (les valeurs
publiées 46.7 et 91.6 ne diffèrent que par l'arrondi des deux moyennes). La
fraction de l'apport solaire est $46.6/173000 = \num{2.7e-4}$, soit
\qty{0.027}{\percent}.
\end{solution}

\begin{solution}{exo:b1:heat-budget:3}
Uranium~: $0.99274\times\num{9.513e-5} + 0.00720\times\num{5.685e-4} =
\num{9.444e-5} + \num{4.09e-6} = \qty{9.85e-5}{W.kg^{-1}}$~; le rare
\ce{^{235}U} en fournit \qty{4}{\percent}. Potassium~: la fraction massique du
\ce{^{40}K} est $\num{1.17e-4}\times39.96/39.10 = \num{1.196e-4}$, donc
$\num{1.196e-4}\times\num{2.847e-5} = \qty{3.40e-9}{W.kg^{-1}}$.
\end{solution}

\begin{solution}{exo:b1:heat-budget:4}
$\qty{1}{\micro cal.cm^{-2}.s^{-1}} = \num{4.184e-6}\,\unit{J.s^{-1}}/\qty{e-4}{m^2}
= \qty{41.84}{mW.m^{-2}}$~: est des États-Unis \qty{33.5}{mW.m^{-2}}, Basin and
Range \qty{58.6}{mW.m^{-2}}, Sierra Nevada \qty{16.7}{mW.m^{-2}}. Avec $A_0D$ en
\unit{mW.m^{-2}} quand $A_0$ est en \unit{\micro W.m^{-3}} et $D$ en \unit{km}~:
$33.5 + 3.0\times7.5 = \qty{56.0}{mW.m^{-2}}$, $58.6 + 3.0\times9.4 =
\qty{86.8}{mW.m^{-2}}$, $16.7 + 3.0\times10 = \qty{46.7}{mW.m^{-2}}$. Le même
granite donne des flux presque deux fois plus forts dans le Basin and Range que
dans la Sierra Nevada~: la différence est dans le \omterm{def:b1:layered-earth:layers}{manteau} en dessous.
\end{solution}

\begin{solution}{exo:b1:heat-budget:5}
$T(10\,\text{km}) = 10 + \frac{0.060\times10^4}{2.5} - \frac{\num{2e-6}\times10^8}{2\times2.5}
= 10 + 240 - 40 = \qty{210}{\degreeCelsius}$. Below, $q_m = 60 - 2.0\times10 =
\qty{40}{mW.m^{-2}}$, un gradient de \qty{16}{K.km^{-1}}, et
$T(35\,\text{km}) = 210 + 16\times25 = \qty{610}{\degreeCelsius}$.
\end{solution}

\begin{solution}{exo:b1:heat-budget:6}
Argilite~: gradient $12.0/0.300 = \qty{40}{K.km^{-1}}$, $q = 1.5\times0.040 =
\qty{60}{mW.m^{-2}}$. Granite~: $10.0/0.500 = \qty{20}{K.km^{-1}}$,
$q = 3.0\times0.020 = \qty{60}{mW.m^{-2}}$. $R(300) = 300/1.5 =
\qty{200}{m^2.K.W^{-1}}$, $R(800) = 200 + 500/3.0 = \qty{366.7}{m^2.K.W^{-1}}$.
Les trois points $(0, 12.0)$, $(200, 24.0)$, $(366.7, 34.0)$ sont alignés, de
pente $22.0/366.7 = \qty{0.0600}{W.m^{-2}}$~: le flux est le même,
\qty{60}{mW.m^{-2}}, dans les deux roches~; seul le gradient change, en
proportion inverse de la conductivité.
\end{solution}

\begin{solution}{exo:b1:heat-budget:7}
$\tau = L^2/\kappa$~: \qty{1}{m}, \qty{e6}{s} ou \num{12}~jours~; \qty{10}{m},
\qty{e8}{s} ou \num{3.2}~ans~; \qty{100}{km}, \qty{e16}{s} ou
\qty{317}{\Ma}~; \qty{2900}{km}, \qty{8.4e18}{s} ou \qty{267}{\Ga}. À
\qty{10}{m} le cycle journalier ne se ressent pas du tout, et le cycle annuel,
plus court que les trois ans nécessaires pour traverser \qty{10}{m}, très
faiblement seulement~: seules les variations qui durent des années ou plus
passent. Le dernier temps vaut $267/4.54 = 59$ fois l'âge de la Terre.
\end{solution}

\begin{solution}{exo:b1:heat-budget:8}
(a) $1646 - 0.54\times50 = \qty{1619}{K}$, soit \qty{1346}{\degreeCelsius},
la valeur publiée de \qty{1350}{\degreeCelsius} arrondie. (b) Gradient moyen
$(0.54 + 0.36)/2 = \qty{0.45}{K.km^{-1}}$, donc $1646 + 0.45\times360 =
\qty{1808}{K}$, \qty{9}{K} (\qty{0.5}{\percent}) au-dessus de \qty{1799}{K}~: le
gradient ne décroît pas linéairement, mais l'estimation est proche.
\end{solution}

\begin{solution}{exo:b1:heat-budget:9}
(a) $G^2 = T_0^2/(\pi\kappa t)$ donne $t = T_0^2/(\pi\kappa G^2)$.
(b) $G = 0.02\times\tfrac59\,\unit{K}/\qty{0.3048}{m} = \qty{0.03645}{K.m^{-1}} =
\qty{36.45}{K.km^{-1}}$; $t = 3900^2/(\pi\times\num{1.2e-6}\times0.03645^2) =
\num{1.521e7}/\num{5.009e-9} = \qty{3.04e15}{s}$, soit \qty{96}{\Ma},
le chiffre de Kelvin. (c) $t = 1300^2/(\pi\times\num{e-6}\times0.036^2) =
\qty{4.15e14}{s} = \qty{13}{\Ma}$~: une température initiale plus réaliste rend la
Terre conductrice plus jeune encore.
\end{solution}

\begin{solution}{exo:b1:heat-budget:10}
(a) Après l'érosion, la nouvelle surface est l'ancienne profondeur $e$, où la
production de chaleur était $A_0' = A_0e^{-e/D}$ et le flux $q(e) = q_m +
A_0De^{-e/D} = q_m + A_0'D$~; mesurée depuis la nouvelle surface, la production
décroît encore comme $A_0'e^{-z'/D}$. Le pluton se déplace le long de la même droite.
(b) La couche a maintenant l'épaisseur $h - e$ avec le même $A$, donc $q_0' = q_m + A(h - e)$~:
sur une droite de pente $h - e$ passant par $q_m$. (c) Seul le modèle exponentiel~: dans le
modèle à couche uniforme, des plutons érodés de profondeurs différentes se
placeraient sur des droites de pentes différentes, et les points d'une province
seraient dispersés. Qu'on observe une seule droite plaide pour une production
décroissant à peu près exponentiellement avec la profondeur.
\end{solution}

\begin{solution}{exo:b1:heat-budget:11}
(a) $\tau_P = RL/(3\kappa)$~: pour \qty{50}{km}, $\num{6.4e6}\times\num{5e4}/\num{3e-6}
= \qty{1.07e17}{s} = \qty{3.4}{\Ga}$~; pour \qty{100}{km}, \qty{6.8}{\Ga}.
(b) $q = kT_i/L = 3\times1300/\num{5e4} = \qty{78}{mW.m^{-2}}$ au départ,
et $78\,e^{-2/3.38} = \qty{43}{mW.m^{-2}}$ après \qty{2}{\Ga}. (c)
$3\times1300/\sqrt{\pi\times\num{e-6}\times\num{6.31e16}} = 3900/\num{4.45e5} =
\qty{8.8}{mW.m^{-2}}$~: la Terre bien brassée perd encore cinq fois plus de chaleur
après deux milliards d'années. Tout son intérieur alimente la surface, là où seule
une enveloppe externe le fait dans une Terre conductrice.
\end{solution}

\begin{solution}{exo:b1:heat-budget:12}
(a) $\sqrt{\pi\kappa\times\qty{1}{\Ma}} = \sqrt{\pi\times\num{e-6}\times\num{3.156e13}}
= \qty{9957}{m}$, donc $C = 3\times1336/9957 = \qty{0.403}{W.m^{-2}}$ pour
$t$ en Ma, \qty{403}{mW.m^{-2}.\Ma^{1/2}}~: \qty{82}{\percent} des 490
ajustés. Le modèle demande un $k\,\Delta T$ plus grand (environ \qty{4900}{W.m^{-1}})~: une
conductivité plus élevée du \omterm{def:b1:layered-earth:layers}{manteau}, ou un intérieur plus chaud, comme le long de l'adiabate.
(b) $490/\sqrt1 = 490$, $490/\sqrt{10} = 155$ et $490/\sqrt{100} =
\qty{49}{mW.m^{-2}}$. (c) $z = 2.326\sqrt{\num{e-6}\times\num{3.156e15}} =
2.326\times\qty{56180}{m} = \qty{131}{km}$.
\end{solution}

\begin{solution}{pb:b1:heat-budget:1}
\textbf{1.} Gradient $= (T_{500} - T_{100})/\qty{0.4}{km}$~: A \qty{15.25}{K.km^{-1}},
B \num{19.50}, C \num{19.50}, D \num{25.50}, E \qty{28.25}{K.km^{-1}}.

\textbf{2.} Le forage perturbe la température de la roche (le fluide de forage
refroidit ou réchauffe les parois), et la perturbation met des semaines à des mois
à s'éteindre. Les dizaines de mètres supérieurs suivent le cycle saisonnier et les
changements récents de la température de surface.

\textbf{3.} $q_0 = k\times$gradient~: A $2.9\times15.25 = \qty{44.225}{mW.m^{-2}}$,
B \qty{50.70}{mW.m^{-2}}, C \qty{60.45}{mW.m^{-2}}, D \qty{68.85}{mW.m^{-2}},
E \qty{84.75}{mW.m^{-2}}.

\textbf{4.} A à D sont sous la moyenne continentale, E au-dessus de
\qty{13.9}{mW.m^{-2}}, soit \qty{20}{\percent}.

\textbf{5.} Le flux est la conductivité fois le gradient~: C conduit mieux (3.1
contre 2.6), donc le même gradient porte \qty{19}{\percent} de chaleur en plus.

\textbf{6.} Les isothermes ne suivent le relief que partiellement~: sous un sommet
elles montent moins que la surface, elles sont donc plus écartées en dessous, et le
gradient mesuré est diminué (augmenté sous une vallée). Une correction de relief est
nécessaire avant de comparer les plutons.

\textbf{7.} $\rho C_U H_U$~: A $2650\times\num{2.0e-6}\times\num{9.85e-5} =
\qty{0.522}{\micro W.m^{-3}}$, B 0.914, C 1.305, D 1.827, E 2.610.

\textbf{8.} Thorium~: A 0.558, B 0.976, C 1.255, D 1.742, E 2.439.

\textbf{9.} Potassium~: A $2650\times0.030\times\num{3.40e-9} = 0.270$, B 0.315,
C 0.360, D 0.378, E 0.405.

\textbf{10.} $A_0$ = 1.35, 2.20, 2.92, 3.95 et \qty{5.46}{\micro W.m^{-3}}.
Le potassium contribue le moins, de \qty{20}{\percent} du total en A jusqu'à
\qty{7}{\percent} en E~; l'uranium et le thorium se partagent le reste presque à parts égales.

\textbf{11.} Oui~: les cinq points montent régulièrement, à moins de
\qty{1.5}{mW.m^{-2}} d'une droite.

\textbf{12.} $\bar A_0 = 15.88/5 = 3.176$, $\bar q_0 = 308.975/5 = 61.795$;
$\sum(A_0 - \bar A_0)^2 = 10.1681$; $\sum(A_0 - \bar A_0)(q_0 - \bar q_0) =
101.146$.

\textbf{13.} $D = 101.146/10.1681 = 9.95$ (en \unit{mW.m^{-2}} par
\unit{\micro W.m^{-3}}, c'est-à-dire en km)~: $D = \qty{9.95}{km}$, et $q_m = 61.795 -
9.95\times3.176 = \qty{30.20}{mW.m^{-2}}$.

\textbf{14.} La pente, environ \qty{10}{km}, est celle des trois provinces~; le
flux réduit, \qty{30}{mW.m^{-2}}, est proche de celui de l'est des États-Unis,
\qty{33}{mW.m^{-2}}~: une province calme avec un \omterm{def:b1:layered-earth:layers}{manteau} froid en dessous, à la
différence du Basin and Range.

\textbf{15.} $q_m/k = \qty{12}{K.km^{-1}}$ et $A_0D^2/k = \num{5.46e-6}\times\num{e8}/2.5
= \qty{218.4}{K}$~: $T(z) = 10 + 12z + 218.4\,(1 - e^{-z/10})$, $z$ en km, $T$ en
\unit{\degreeCelsius}.

\textbf{16.} $q(z) = 30 + 54.6\,e^{-z/10}$~: \qty{50.1}{mW.m^{-2}} à \qty{10}{km},
\num{37.4} à \qty{20}{km} et \qty{31.7}{mW.m^{-2}} à \qty{35}{km}.

\textbf{17.} $30/84.6 = \qty{35}{\percent}$~: près des deux tiers du flux
au-dessus de E sont produits dans sa couche radiogénique.

\textbf{18.} $T(10) = 10 + 120 + 218.4\times0.632 = \qty{268}{\degreeCelsius}$.

\textbf{19.} $T(20) = 10 + 240 + 218.4\times0.865 = \qty{439}{\degreeCelsius}$.

\textbf{20.} Sous A, $A_0D^2/k = 1.35\times100/2.5 = \qty{54.0}{K}$~:
$T(35) = 10 + 420 + 54.0\times0.970 = \qty{482}{\degreeCelsius}$.

\textbf{21.} De $A_0D^2/k$~: \qty{54}{K} sous A et \qty{218}{K} sous E. Le flux
profond est le même~; la \omterm{def:b1:layered-earth:layers}{croûte} de E est quatre fois plus riche en sources de chaleur, et la
chaleur qu'elles produisent doit être conduite vers le haut à travers la même roche, ce qui demande un
gradient plus raide sur toute la hauteur.

\textbf{22.} $1300 = 10 + 218.4 + 12z$ donne $z = \qty{89.3}{km}$ sous E~;
$1300 = 10 + 54.0 + 12z$ donne \qty{103.0}{km} sous A.

\textbf{23.} Les deux se situent dans l'intervalle de la \omterm{def:b1:layered-earth:lithosphere}{lithosphère} continentale,
d'une centaine de kilomètres ou moins sous les régions jeunes à environ \qty{250}{km} sous les
cratons, du côté mince~: avec un flux réduit de \qty{30}{mW.m^{-2}} et cette
conductivité, le \omterm{def:b1:heat-budget:geotherm}{géotherme} conductif atteint les températures du \omterm{def:b1:layered-earth:layers}{manteau} vers
\qty{100}{km}~; sous le pluton le plus radioactif, \qty{14}{km} plus tôt.

\textbf{24.} $T(35) = 10 + 420 + 218.4\times(1 - e^{-3.5}) = 430 + 211.8$~: le
\omterm{def:b1:layered-earth:discontinuity}{Moho} sous le pluton E est à
\textbf{\qty[text-series-to-math=true]{642}{\degreeCelsius}}.
\end{solution}
